% ============================================================
% 第 7 章 其他常用变压器（教材第二篇，印刷页 68--75）
% 节名：第一节 自耦变压器；第二节 三绕组变压器；第三节 互感器
% ============================================================
\chapter{其他常用变压器}

前面各章讨论的变压器，每相只有一个一次绕组和一个二次绕组，称为\textbf{双绕组变压器}——它是最普通、最常用的变压器。除此之外还有许多型式和用途的变压器，本章介绍较常用的自耦变压器、三绕组变压器和互感器。它们的基本原理仍是双绕组变压器那一套，特殊之处在于\textbf{绕组的连接方式}或\textbf{绕组的数目}改变了功率传递的路径。

\section{自耦变压器}

\subsection{结构特点和变比}
普通双绕组变压器的一、二次绕组之间相互绝缘，只有磁的耦合，没有直接电的联系。如果把普通双绕组变压器的一次绕组和二次绕组\textbf{串联}起来作为新的一次绕组，而原二次绕组还同时作为新的二次绕组，便得到一台降压自耦变压器。

设接成普通双绕组时一、二次绕组匝数分别为 $N_1$、$N_2$；改为自耦连接时，一次绕组匝数 $N_{1a}=N_1+N_2$，二次绕组匝数 $N_{2a}=N_2$。当绕组 $Ax$ 两端外加交变正弦电压 $U_{1a}$ 时，铁芯内产生交变磁通，在 $Ax$ 间感应电动势
\begin{equation}
E_{1a}=4.44f\left(N_1+N_2\right)\Phi_m=E_1+E_2 ,
\end{equation}
在二次绕组 $ax$ 间感应电动势 $E_{2a}=4.44fN_2\Phi_m=E_2$。忽略漏阻抗压降，自耦变压器的变比为
\begin{equation}
k_a=\frac{N_{1a}}{N_{2a}}=\frac{N_1+N_2}{N_2}=1+\frac{N_1}{N_2}=1+\frac{1}{k},
\end{equation}
其中 $k=N_1/N_2$ 为接成双绕组变压器时的变比。

可见自耦变压器的变压原理与双绕组变压器相同，只是由于一、二次绕组有直接的电联系、一次绕组包含了二次绕组，因此自耦变压器的变比 $k_a$ 比接成双绕组时的变比 $k$ 大。用自耦变压器升压时，可把电源加在匝数为 $N_2$ 的绕组上，匝数为 $N_{1a}$ 的绕组感应电动势就比外加电压高。

\subsection{基本电磁关系}
\subsubsection{电压、电流关系}
对降压自耦变压器，二次绕组（$ax$ 段）的电流为
\begin{equation}
\dot I_{2}=\dot I_1+\dot I_{2a},
\end{equation}
即在忽略励磁电流的情况下，由磁平衡关系 $N_1\dot I_1+N_2\dot I_2=0$（对串联段而言），可得公共绕组中的电流为两电流之差。故自耦变压器二次电流为"一次电流与公共段电流的相量和"，公共绕组段的电流较小——这正是自耦变压器省材料的根本原因。

\subsubsection{功率关系}
作为普通双绕组变压器时，额定容量为
\begin{equation}
S_N=U_{1N}I_{1N}=U_{2N}I_{2N}.
\end{equation}
接成自耦变压器时，额定容量为
\begin{equation}
S_{aN}=U_{1aN}I_{1aN}=S_N+\frac{U_{2N}}{U_{1N}}S_N=S_N+S_Q ,
\end{equation}
即
\begin{equation}
S_{aN}=S_N+S_Q .
\end{equation}

\keyformula{S_{aN}=S_N+S_Q,\qquad S_N=\left(1-\frac{1}{k_a}\right)S_{aN}}

自耦变压器的视在功率由两部分组成：
\begin{itemize}
  \item \textbf{感应功率} $S_N$：借助于 $Aa$ 段绕组和 $ax$ 段绕组之间的电磁感应关系，从一次侧传递到二次侧的功率。这部分功率与普通双绕组变压器相同，变压器的\textbf{主要尺寸和材料消耗由这部分功率决定}，所以 $S_N$ 称为\textbf{计算容量}。
  \item \textbf{传导功率} $S_Q$：通过电路的直接联系，从一次绕组直接传递到负载侧的功率。它\textbf{不需要增加变压器的计算容量}。
\end{itemize}

\subsection{自耦变压器的优缺点}
\begin{enumerate}
  \item \textbf{体积小、质量轻、价格低}，便于运输和安装。自耦变压器的计算容量 $S_N=(1-1/k_a)S_{aN}<S_{aN}$，故在同样额定容量下，主要尺寸和所用铜线、硅钢片及其他结构材料都比较少；铜线和硅钢片减少，铜耗和铁耗也相应减少，效率提高。$k_a$ 越接近 1，计算容量占额定容量的比例越小，这些优点越显著。因此自耦变压器适合于一、二次电压相差不大的场合，一般 $k_a<2$。
  \item 将低压侧短路、从高压侧看，其\textbf{短路阻抗与普通双绕组变压器相同}，但由于自耦变压器的电压基值较高，故\textbf{短路阻抗标幺值比普通双绕组变压器小，短路电流较大}。
  \item 高、低压绕组之间具有\textbf{直接的电联系}：一方面要求低压绕组具有与高压绕组相同的绝缘水平，另一方面使\textbf{过电压保护变得复杂}。
\end{enumerate}

\begin{center}
\begin{tabular}{@{}p{2.6cm}p{5.0cm}p{5.0cm}@{}}
\toprule
项目 & 双绕组变压器 & 自耦变压器 \\
\midrule
电的联系 & 只有磁耦合，无直接电联系 & 有直接电联系（一次含二次） \\
变比 & $k=N_1/N_2$ & $k_a=1+1/k=N_{1a}/N_{2a}$ \\
额定容量 & $S_N=U_{1N}I_{1N}$ & $S_{aN}=S_N+S_Q$ \\
计算容量 & 等于额定容量 & $(1-1/k_a)S_{aN}<S_{aN}$ \\
短路电流 & 较小 & 标幺短路阻抗小，短路电流较大 \\
绝缘与保护 & 简单 & 低压侧需按高压绝缘；过电压保护复杂 \\
适用场合 & 通用 & $k_a<2$，电压相近的场合 \\
\bottomrule
\end{tabular}
\end{center}

\corepoint{自耦变压器的全部"异常"都源于一句话：\textbf{一次绕组包含了二次绕组}。由此推出三件事——变比 $k_a=1+1/k$ 而不是 $k$；容量由"感应功率 + 传导功率"两部分构成，故计算容量小、省材料；高低压直接电联系，故低压侧要按高压绝缘、短路电流大、过电压保护复杂。答题时先写这句，后面自然出来。}

\section{三绕组变压器}

三绕组变压器有高压、中压和低压三个绕组，大多用于二次侧需要两种不同电压等级的电力系统。三绕组变压器相当于不同变比的两台双绕组变压器，但使用三绕组变压器比用两台更经济。

\subsection{结构特点和容量分配}
三个绕组一般均套装在同一铁芯柱上。为便于绝缘，\textbf{高压绕组通常放置于最外层}。对于\textbf{降压}变压器，中压绕组放在中间层，低压绕组放在最里层；对于\textbf{升压}变压器，为使漏磁场分布均匀、漏抗分配合理，不致因低压和高压绕组相距较远而造成漏磁通及其附加损耗增加，保证有较好的电压变化率和运行性能，通常将中压绕组放在最里层、低压绕组放在中间层。

GB 1094.1—1996 规定，三相三绕组变压器的标准联结组有 \textbf{YNyn0d11} 和 \textbf{YNyn0y0} 两种。

三个绕组的容量可以相同也可以不同。三绕组变压器的\textbf{额定容量是指三个绕组中容量最大那个绕组的容量}。若以额定容量为 100\%，则三个绕组容量的配合有 100/100/50、100/50/100 和 100/100/100 三种。

\subsection{基本方程式和等效电路}
设绕组 1 是一次绕组，另外两个是二次绕组，则三绕组变压器的变比分别为
\begin{equation}
k_{12}=\frac{N_1}{N_2},\qquad k_{13}=\frac{N_1}{N_3}.
\end{equation}

三绕组变压器在磁路方面有三个绕组相互耦合，它们产生的磁通可分为三种：
\begin{itemize}
  \item \textbf{自漏磁通}：只与一个绕组相交链；
  \item \textbf{互漏磁通}：与两个绕组交链而不与第三个绕组相交链；
  \item 以上两部分漏磁通通过油或空气闭合，\textbf{不受铁芯磁路饱和的影响}；
  \item \textbf{主磁通}：与三个绕组同时交链，通过铁芯磁路闭合，\textbf{受铁芯饱和影响}。
\end{itemize}

主磁通由三个绕组的合成磁动势产生，故励磁磁动势为
\begin{equation}
N_1\dot I_m=N_1\dot I_1+N_2\dot I_2+N_3\dot I_3
\quad\Longrightarrow\quad
\dot I_1+\dot I_2'+\dot I_3'=\dot I_m ,
\end{equation}
其中 $\dot I_2'$、$\dot I_3'$ 分别表示归算到一次侧的两个二次绕组电流。忽略励磁电流时
\begin{equation}
\dot I_1+\dot I_2'+\dot I_3'=0 .
\end{equation}

由于存在互漏磁通，每个绕组内的漏磁压降都受另外两个绕组的影响。若按双绕组变压器的分析方法建立等效电路，三绕组变压器每个绕组的电抗必须是一个既包含自感电抗、又包含互感电抗的\textbf{等效值}。用 $X_1$、$X_2'$、$X_3'$ 和 $Z_1$、$Z_2'$、$Z_3'$ 表示三个绕组的等效电抗和等效阻抗，可列出电压方程
\begin{equation}
\dot U_1=-\dot E_1+\dot I_1Z_1,\qquad
\dot U_2'=\dot E_1-\dot I_2'Z_2',\qquad
\dot U_3'=\dot E_1-\dot I_3'Z_3' .
\end{equation}
据此可得三绕组变压器的\textbf{简化等效电路}：三个等效阻抗呈星形连接，公共点接励磁支路。

\subsection{参数测定}
简化等效电路中的 $Z_1$、$Z_2'$、$Z_3'$ 可通过\textbf{三次短路试验}确定：
\begin{enumerate}
  \item 绕组 2 短路、绕组 3 开路，在绕组 1 上加低电压做短路试验，测得 1、2 之间的短路阻抗 $Z_{k12}=Z_1+Z_2'$；
  \item 绕组 3 短路、绕组 2 开路，测得 $Z_{k13}=Z_1+Z_3'$；
  \item 绕组 3 短路、绕组 1 开路，测得 $Z_{k23}=Z_2'+Z_3'$。
\end{enumerate}
联解三式得
\begin{equation}
Z_1=\tfrac12\left(Z_{k12}+Z_{k13}-Z_{k23}'\right),\quad
Z_2'=\tfrac12\left(Z_{k12}+Z_{k23}'-Z_{k13}\right),\quad
Z_3'=\tfrac12\left(Z_{k13}+Z_{k23}'-Z_{k12}\right).
\end{equation}
其中 $Z_1$、$Z_2'$、$Z_3'$ 均为等效阻抗，其大小与三个绕组之间的布置有关；在某种特定布置下，其中某个等效电抗可能变成负值，此时该电抗可当作\textbf{容抗}计算。

\corepoint{三绕组变压器的等效电路不是"三个绕组各串一个阻抗再并到励磁支路"，而是\textbf{三个等效阻抗星形连接后接励磁支路}——因为互漏磁通使每个绕组的漏抗不再是独立的物理量，只能由三次短路试验"解方程"求出。若算出某个 $X$ 为负值不要慌，那是互漏磁通造成的等效结果，按容抗处理即可。}

\section{互感器}

互感器是测量用设备，分\textbf{电压互感器}和\textbf{电流互感器}两种，主要在高电压、大电流测量以及各种继电保护中使用。互感器的工作原理与普通变压器相同，但运行状态与使用禁忌截然相反——这是本节的核心。

\subsection{电压互感器}
电压互感器一次绕组匝数 $N_1$ 较多，直接接于被测的高压电路上；二次绕组匝数较少，一般只有一匝或几匝，与电压表或其他仪表的电压线圈相接。由于电压表或仪表的电压线圈内阻抗很大，所以电压互感器工作时相当于变压器的\textbf{空载运行}。

电压互感器二次绕组的额定电压通常取 100\,V。忽略漏阻抗压降，则
\begin{equation}
\frac{U_1}{U_2}=\frac{N_1}{N_2}=k_u\quad\Longrightarrow\quad U_1=k_uU_2 ,
\end{equation}
因此利用一、二次绕组不同的匝数比可将电路上的高电压变为低电压来测量。

\subsubsection{误差与准确级}
电压互感器的准确度由\textbf{变比误差}和\textbf{相位误差}衡量：
\begin{itemize}
  \item 变比误差：一次电压 $U_1$ 与二次电压归算值 $k_uU_2$ 的算术差的相对值，即 $\left[(U_1-k_uU_2)/U_1\right]\times100\%$；
  \item 相角误差：$\dot U_1$ 与 $-k_u\dot U_2$ 之间的相位差。
\end{itemize}
从变压器相量图可以看出，\textbf{励磁电流和一、二次绕组的漏阻抗压降}是产生这两种误差的主要原因。因此在设计电压互感器时，大都采用高磁导率的硅钢片，尽量减小磁路中的气隙，并使铁芯不饱和，一般取磁密为 0.6～0.8\,T。按误差大小，电压互感器的准确度分为 0.1、0.2、0.5、1.0 和 3.0 五级。

\subsubsection{使用注意}
\begin{enumerate}
  \item \textbf{二次绕组不允许短路}，否则会产生很大的短路电流；
  \item 互感器的\textbf{铁芯和二次绕组的一端必须可靠接地}，以保证操作人员的安全。
\end{enumerate}

\subsection{电流互感器}
电流互感器一次绕组匝数 $N_1$ 较少，通常为一匝或几匝，串入需要测量电流的电路中；二次绕组匝数 $N_2$ 较多，与电流表或其他仪表的电流线圈相接。电流表或仪表中电流线圈的阻抗很小，所以电流互感器工作时相当于普通变压器的\textbf{短路运行}。

电流互感器二次绕组的额定电流一般为 5\,A 或 1\,A。忽略励磁电流，根据磁平衡方程式 $I_1N_1+I_2N_2=0$，在数值上可得
\begin{equation}
I_1=\frac{N_2}{N_1}I_2=k_iI_2 ,
\end{equation}
即用匝数比把大电流变成小电流来测量。

\subsubsection{误差与准确级}
与电压互感器类似，电流互感器的误差也分变比误差与相角误差，主要由励磁电流和漏阻抗压降引起。为减小误差，应采用高磁导率材料、减小磁路气隙、降低铁芯磁密（即增大铁芯截面）。按误差大小分为 0.1、0.2、0.5、1.0、3.0 等级。

\subsubsection{使用注意}
\begin{enumerate}
  \item \textbf{二次绕组绝对不允许开路}。二次开路时，一次电流全部变为励磁电流，铁芯中磁通和磁密急剧增大，铁耗剧增、铁芯严重过热；同时在二次绕组两端感应出很高的电压，危及人身与设备安全。
  \item \textbf{铁芯和二次绕组的一端必须可靠接地}。
\end{enumerate}

\begin{center}
\begin{tabular}{@{}p{2.3cm}p{3.3cm}p{3.6cm}p{4.0cm}@{}}
\toprule
项目 & 电压互感器 & 电流互感器 & 记忆要点 \\
\midrule
一次绕组 & 匝数多，并联在被测电路 & 匝数少（1～几匝），串入电路 & 并联/串联 \\
工作状态 & 相当于空载运行 & 相当于短路运行 & 与变压器对应 \\
二次额定值 & 100\,V & 5\,A 或 1\,A & 标准化量程 \\
主要误差源 & 励磁电流、漏阻抗压降 & 励磁电流、漏阻抗压降 & 都靠高磁导率、低磁密减小 \\
使用禁忌 & \textbf{二次不许短路} & \textbf{二次不许开路} & 一字之差，后果都严重 \\
安全措施 & 铁芯与二次一端接地 & 铁芯与二次一端接地 & 防止高压窜入低压 \\
\bottomrule
\end{tabular}
\end{center}

\corepoint{记住这一组"反着来"的关系：电压互感器 = \textbf{空载}运行，因此怕\textbf{短路}（短路即成为短路运行，电流剧增）；电流互感器 = \textbf{短路}运行，因此怕\textbf{开路}（开路即失去去磁作用，铁芯饱和、二次出现高压）。两条安全措施完全相同：铁芯与二次一端必须接地。}
